\chapter{Graphs and Hypergraphs}
\label{ch:graphs}


Section~\ref{sec:graphs} introduced the \B{Graph} expression: how to build one,
read it back, take its adjacency matrix, and ask for a shortest path or a
colouring. This chapter continues from there and covers the rest of the
subsystem. It is about the questions that make graph theory a subject and not
only a data structure. How many spanning trees does a network have? Where does
the traffic concentrate? How much can flow from a source to a sink, and which
edges are the bottleneck? Are two graphs the same graph in disguise? Can a
graph be drawn in the plane without crossings? The second half of the chapter
drops the assumption that an edge joins exactly two things. A \emph{hypergraph}
allows an edge to join any number of vertices, which is the natural model for
co-authorship, chemical reactions, database relations, and the rewriting states
of Wolfram-model physics.

Three themes run through the chapter. The first is the mathematics: most
sections lead with a theorem (Kirchhoff, Menger, Gallai, Kuratowski, Berge) and
then check it at the prompt. The second is exactness. Mathilda gives an exact
answer whenever the problem has one, and a hard search that runs out of budget
leaves the call unevaluated; it never returns a result that is merely plausible.
The third is speed. Each algorithm here was benchmarked against Mathematica~15
and against networkx or xgi. The last section gives the figures, including the
one case where Mathilda is slower.

\section{A laboratory of named graphs}
\label{sec:graphs-families}

Graph theory has its favourite examples, and the one that disproves more
conjectures than any other is the \emph{Petersen graph}.\index{Petersen graph}
It has ten vertices and fifteen edges, and every vertex has degree three.
\B{PetersenGraph} generates it, and a few measurements show how tightly it is
built:

\begin{codepairs}
\pair{graphs-hypergraphs/families/1}{The Petersen graph, printed as a summary
like every graph.}
\pair{graphs-hypergraphs/families/2}{Ten vertices, fifteen edges, and only one
degree occurs: the graph is \emph{3-regular}.}
\pair{graphs-hypergraphs/families/3}{\B{GraphDiameter} and \B{GraphRadius} are
both \(2\): any two vertices are at most two steps apart.
\B{MeanGraphDistance} is exact.}
\pair{graphs-hypergraphs/families/4}{\B{GraphDensity}: \(15\) of the
\(\binom{10}{2} = 45\) possible edges are present.}
\end{codepairs}

The mean distance \(5/3\) is easy to confirm by hand. Each vertex has three
neighbours at distance \(1\) and six more vertices at distance \(2\), so the
mean over its nine partners is \((3 + 12)/9\). Mathilda keeps that as the
rational \(5/3\) and does not round it to \(1.66667\). This happens throughout
the subsystem: a distance or a coefficient over integer data comes back exact.

The Petersen graph is one of a large set of generators. \B{HypercubeGraph},
\B{WheelGraph}, \B{GridGraph}, \B{TuranGraph}, \B{CirculantGraph},
\B{CompleteKaryTree}, and others give the standard families on the vertices
\(1, \dots, n\). \B{GraphDistanceMatrix} gives every pairwise distance at once:

\begin{codepairs}
\pair{graphs-hypergraphs/families/5}{Four standard families: the 3-cube, a
wheel, a \(3\times4\) grid, and the Tur\'an graph \(T(6,3)\), which is
\(K_{2,2,2}\).}
\pair{graphs-hypergraphs/families/6}{All-pairs distances in the 5-cycle. Each
row is a cyclic shift of \(0,1,2,2,1\).}
\end{codepairs}

\calloutnote{Under the hood}{\B{GraphDistanceMatrix} on an unweighted graph runs
a \emph{bit-parallel multi-source breadth-first search} (Then et~al., 2015).
Up to 256 sources advance together, one bit per source at each vertex, so each
sweep over the adjacency structure settles 256 rows of the matrix. The result is a packed
integer matrix. A weighted graph goes to Dijkstra instead, and exact weights
give an exact matrix.}

\section{Drawing graphs}
\label{sec:graphs-drawing}

A graph has no shape of its own. A drawing is a choice of position for every
vertex, and a good choice makes structure visible: symmetry, clusters,
bottlenecks. \B{GraphPlot} makes that choice automatically. Its default for
general graphs is \emph{stress majorization}\index{graph drawing!stress
majorization}, which places the vertices so that the straight-line distance
between every two of them is as close as possible to their graph distance.
For the Petersen graph the result is symmetric but has three crossings:

\plotcell{graphs-hypergraphs/petersen}

The drawing most books use is different: an outer pentagon, an inner pentagram,
and five spokes. \mcode{VertexCoordinates} fixes the positions yourself. Here they
are two rings of five points, listed in \B{VertexList} order:

\plotcell{graphs-hypergraphs/petersen-classic}

Both pictures show the same graph. The first is the one the layout algorithm
finds from the distances alone. The second uses outside knowledge of the
graph's symmetry. Later sections reuse the classic coordinates \mcode{pent}.

The automatic layout adapts to the kind of graph. A grid comes out as a grid,
and the 3-cube as the familiar Necker drawing, with \mcode{VertexLabels} showing
the vertex names:

\plotcell{graphs-hypergraphs/grid}

\plotcell{graphs-hypergraphs/families}

\calloutnote{Under the hood}{\B{GraphPlot} (\texttt{src/graph/glayout.c})
starts the stress layout from a Pivot MDS embedding and refines it by SMACOF
iterations on the breadth-first distance matrix. For components of up to 60
vertices it also tries circular and Tutte starting points, and among the
drawings within \(10\%\) of the lowest stress it keeps the one with the fewest
edge crossings. Trees and directed acyclic graphs are drawn in layers instead,
with dummy vertices on long edges and barycentre ordering to reduce crossings.
Components are packed side by side. The layout uses no random numbers, so the
same graph always gives the same picture.}

The rest of the chapter uses three more options. \mcode{GraphHighlight} draws chosen
vertices and edges in red, \mcode{VertexStyle} colours vertices individually,
and \mcode{EdgeLabels} writes text on edges. Hypergraphs have their own drawing
function, \mcode{HypergraphPlot}, which appears in \S\ref{sec:hypergraphs}.

\section{Directed graphs and dependency}
\label{sec:graphs-directed}

A directed edge \(u \to v\) often means ``\(u\) must happen before \(v\)''. A
build system, a course prerequisite chart, and a spreadsheet's formula
dependencies are all directed graphs of this kind. Getting dressed is a small
example: socks and trousers go on before shoes, the shirt before the tie, and so
on.\index{directed acyclic graph}

\begin{codepairs}
\pair{graphs-hypergraphs/directed/1}{Seven garments and the order constraints
between them.}
\pair{graphs-hypergraphs/directed/2}{\B{TopologicalSort} lists the vertices so
that every edge points forward: a valid order in which to get dressed.}
\pair{graphs-hypergraphs/directed/3}{A topological order exists exactly when the
graph has no directed cycle. \B{AcyclicGraphQ} checks this.}
\pair{graphs-hypergraphs/directed/4}{Add an edge from the jacket back to the
shirt and the constraints become impossible to satisfy.}
\end{codepairs}

\B{GraphPlot} draws a directed acyclic graph in layers, with every arrow
pointing down. A topological order is any left-to-right, top-to-bottom reading
that respects the arrows:

\plotcell{graphs-hypergraphs/dress}


\begin{theorem}
A directed graph has a topological order if and only if it is acyclic.
\end{theorem}

One direction is immediate: along a directed cycle, no vertex can come before
all the others. For the other, every finite acyclic graph has a vertex with no
incoming edges. Put it first, delete it, and repeat. That argument is Kahn's
algorithm, and it runs in time linear in the size of the graph.

When a directed graph does have cycles, its \emph{strongly connected
components}\index{strongly connected component} are the maximal sets of vertices
that can all reach one another. Contracting each component to a single vertex
always leaves an acyclic graph, so every directed graph is a DAG of strongly
connected pieces:

\begin{codepairs}
\pair{graphs-hypergraphs/directed/5}{A graph with a 3-cycle, a 2-cycle, and a
vertex that only points in.}
\pair{graphs-hypergraphs/directed/6}{\B{StronglyConnectedComponents} finds the
three pieces.}
\pair{graphs-hypergraphs/directed/7}{\B{ReverseGraph} flips every edge. The
strongly connected components are unchanged.}
\end{codepairs}

A directed graph also has an \emph{oriented incidence matrix} \(B\), with one row
per vertex and one column per edge: \(-1\) where the edge leaves and \(+1\) where
it arrives. Its product with its own transpose is the Laplacian of the
underlying undirected graph, whatever orientation was chosen:

\begin{codepairs}
\pair{graphs-hypergraphs/directed/8}{\B{IncidenceMatrix} of a directed graph
records each edge's tail and head.}
\pair{graphs-hypergraphs/directed/9}{\(BB^{\mathsf T} = L\) holds exactly.}
\end{codepairs}

\calloutnote{Theory}{The identity \(L = BB^{\mathsf T}\) shows at once that the
Laplacian is positive semidefinite, since
\(x^{\mathsf T} L x = \lVert B^{\mathsf T} x\rVert^2 = \sum_{u\to v}(x_u - x_v)^2
\ge 0\). The quadratic form measures how much a vertex function \(x\) varies
across the edges, and that is why the small Laplacian eigenvalues describe the
smooth, slowly varying structure of a graph: its clusters and bottlenecks.}

\calloutnote{Pitfall}{\B{ConnectedComponents} on a \emph{directed} graph gives
the components of the underlying undirected graph in Mathilda, which is the same
as \B{WeaklyConnectedComponents}. Mathematica returns the strongly connected
components here. Use \B{StronglyConnectedComponents} when you mean mutual
reachability.}

\section{The Laplacian and counting spanning trees}
\label{sec:graphs-laplacian}

Section~\ref{ssec:graph-matrix} used the adjacency matrix \(A\). The matrix that
turns out to be more useful is the \emph{Laplacian} or \emph{Kirchhoff matrix}
\(L = D - A\), where \(D\) is the diagonal matrix of vertex
degrees.\index{Laplacian matrix}\index{Kirchhoff matrix} \B{KirchhoffMatrix}
builds it. Every row of \(L\) sums to zero, so \(L\) is always singular, and the
multiplicity of the eigenvalue \(0\) is the number of connected components.

\begin{codepairs}
\pair{graphs-hypergraphs/spectra/1}{Name the Petersen graph once for this
session.}
\pair{graphs-hypergraphs/spectra/2}{The Laplacian of the 4-cycle has degrees on
the diagonal and \(-1\) at each edge.}
\pair{graphs-hypergraphs/spectra/3}{The Petersen Laplacian has spectrum
\(\{5^4, 2^5, 0^1\}\). There is one zero because the graph is connected.}
\pair{graphs-hypergraphs/spectra/4}{Delete one row and the matching column and
take the determinant. The result, \(2000\), is the number of spanning trees.}
\pair{graphs-hypergraphs/spectra/5}{The adjacency spectrum is
\(3^1, 1^5, (-2)^4\). It is related to the Laplacian spectrum by
\(\lambda_L = 3 - \lambda_A\).}
\end{codepairs}

\begin{theorem}[Kirchhoff's matrix-tree theorem]
\index{matrix-tree theorem}
Let \(G\) be a graph on \(n\) vertices with Laplacian \(L\). Every cofactor of
\(L\) equals the number \(\tau(G)\) of spanning trees of \(G\). Equivalently, if
\(\mu_1, \dots, \mu_{n-1}\) are the nonzero Laplacian eigenvalues,
\[
  \tau(G) \;=\; \frac{1}{n}\prod_{i=1}^{n-1} \mu_i .
\]
\end{theorem}

The session verifies the theorem both ways. The cofactor is \(2000\), and the
spectrum gives \(\tfrac{1}{10}\cdot 5^4 \cdot 2^5 = 2000\) as well. Because the
Laplacian of an unweighted graph is an integer matrix, \B{Det} computes the
cofactor exactly, with no rounding. Counting the spanning trees of a
large graph therefore costs one exact determinant, however many trees there
are.

\calloutnote{Theory}{The factor \(3 - \lambda_A\) is specific to regular graphs:
when every degree is \(k\), \(D = kI\) and \(L = kI - A\) has the same
eigenvectors as \(A\). For irregular graphs the two spectra carry different
information, and the \emph{normalised} Laplacian \(D^{-1/2} L D^{-1/2}\) is often
the better object.}

\section{Trees and Cayley's formula}
\label{sec:graphs-trees}

A \emph{tree}\index{tree} is a connected graph with no cycles. Equivalently, it
is a connected graph on \(n\) vertices with exactly \(n - 1\) edges, or a graph
in which every two vertices are joined by exactly one path. Trees are the
skeletons of graphs: every connected graph contains a \emph{spanning tree} that
reaches all its vertices, and \B{FindSpanningTree} returns one.

\begin{codepairs}
\pair{graphs-hypergraphs/trees/3}{\B{CompleteKaryTree}\mcode{[3, 2]} is the
complete binary tree with three levels.}
\pair{graphs-hypergraphs/trees/4}{Seven vertices and six edges: one fewer edge
than vertices, as for every tree.}
\pair{graphs-hypergraphs/trees/5}{A spanning tree of the Petersen graph keeps
all ten vertices and nine of the fifteen edges.}
\pair{graphs-hypergraphs/trees/6}{\B{TreeGraphQ} confirms it is a tree.}
\end{codepairs}

Trees are drawn in layers from a root. A complete ternary tree with four levels
has \(1 + 3 + 9 + 27 = 40\) vertices:

\plotcell{graphs-hypergraphs/tree}

A spanning tree of the Petersen graph, drawn in red on the classic layout, uses
nine of the fifteen edges and reaches every vertex:

\plotcell{graphs-hypergraphs/spanning}


How many spanning trees does the complete graph \(K_n\) have? The matrix-tree
theorem of \S\ref{sec:graphs-laplacian} answers the question for each \(n\)
separately, and the answers follow a striking pattern:

\begin{codepairs}
\pair{graphs-hypergraphs/trees/1}{The number of spanning trees of \(K_n\) for
\(n = 2, \dots, 8\), each an exact cofactor of the Laplacian.}
\pair{graphs-hypergraphs/trees/2}{The same numbers are \(n^{n-2}\).}
\end{codepairs}

\begin{theorem}[Cayley's formula]
\index{Cayley's formula}
There are exactly \(n^{n-2}\) trees on \(n\) labelled vertices.
\end{theorem}

The spanning trees of \(K_n\) are exactly the labelled trees on its vertices, so
the two computations count the same thing. The Laplacian of \(K_n\) is
\(nI - J\), where \(J\) is the all-ones matrix. Its eigenvalues are \(0\) once
and \(n\) with multiplicity \(n - 1\), and the spectral form of the matrix-tree
theorem gives \(\tfrac1n\, n^{n-1} = n^{n-2}\). The session checked the formula
for seven values of \(n\). The eigenvalue argument proves it for all of them.

\calloutnote{Pitfall}{In this version \B{FindSpanningTree} ignores
\B{EdgeWeight}: on a weighted graph it returns \emph{a} spanning tree, not
necessarily a \emph{minimum} one. Mathematica returns a minimum spanning tree.
Until this is fixed, do not rely on the total weight of the tree it
returns.}

\section{Editing and combining graphs}
\label{sec:graphs-editing}

Graphs are immutable expressions, so an edit returns a new graph and leaves the
original unchanged. \B{VertexDelete}, \B{EdgeAdd}, \B{VertexAdd}, and
\B{EdgeDelete} make local changes. \B{GraphComplement}, \B{LineGraph},
\B{ReverseGraph}, and \B{UndirectedGraph} transform a whole graph.
\B{GraphUnion}, \B{GraphIntersection}, and \B{GraphDifference} combine two.

\begin{codepairs}
\pair{graphs-hypergraphs/editing/1}{Deleting a vertex also deletes its incident
edges. The 5-cycle becomes a path on four vertices.}
\pair{graphs-hypergraphs/editing/2}{Adding the closing edge turns a path into a
triangle.}
\pair{graphs-hypergraphs/editing/3}{The complement of \(C_5\) joins exactly the
pairs that \(C_5\) does not\ldots}
\pair{graphs-hypergraphs/editing/4}{\ldots and that complement is again a
5-cycle. \(C_5\) is \emph{self-complementary}.}
\pair{graphs-hypergraphs/editing/5}{\B{LineGraph} has a vertex for each edge. The
three spokes of a star all meet at the hub, so they form a triangle.}
\pair{graphs-hypergraphs/editing/6}{\B{GraphUnion} merges the vertex and edge
sets. The shared vertex \(3\) glues the two paths together.}
\end{codepairs}

\calloutnote{Performance}{A graph built by \B{Graph} is validated once, and the
result is kept in a small \emph{validated-graph memo} keyed by the expression
node: the vertex hash index, the integer-endpoint edge arrays, and the
compressed adjacency lists. Every later call on the same object starts from
those integer arrays, so predicates such as \B{SimpleGraphQ} take \(O(1)\) time
and editing operations never re-hash vertices. That is why deleting \(1000\)
vertices from a \(10^5\)-vertex graph takes about \(4\) ms in Mathilda and
\(52\) ms in Mathematica.}

\section{Tours: Euler and Hamilton}
\label{sec:graphs-tours}

Two old problems ask for a closed walk that uses every element of a graph once.
An \emph{Eulerian cycle}\index{Eulerian cycle} uses every \emph{edge} exactly
once. A \emph{Hamiltonian cycle}\index{Hamiltonian cycle} visits every
\emph{vertex} exactly once. The two problems sound similar and are very
different in difficulty.

\begin{theorem}[Euler, 1736]
A connected graph has an Eulerian cycle if and only if every vertex has even
degree.
\end{theorem}

That condition is checked in linear time, and Hierholzer's algorithm constructs
the cycle in linear time too. \B{EulerianGraphQ} and \B{FindEulerianCycle} do
exactly this:

\begin{codepairs}
\pair{graphs-hypergraphs/tours/1}{In \(K_5\) every degree is \(4\), which is even.
In \(K_4\) every degree is \(3\), which is odd.}
\pair{graphs-hypergraphs/tours/2}{An Euler tour of \(K_5\): all ten edges, each
traversed once, ending where it began.}
\end{codepairs}

No comparable condition is known for Hamiltonian cycles. Deciding whether one
exists is NP-complete, so \B{FindHamiltonianCycle} has to search. Three graphs
show the three possible outcomes:

\begin{codepairs}
\pair{graphs-hypergraphs/tours/3}{The 3-cube has a Hamiltonian cycle, a
\emph{Gray code} through the eight vertices.}
\pair{graphs-hypergraphs/tours/4}{The Petersen graph has none. The empty list is
a proof of non-existence, not a failure to search.}
\pair{graphs-hypergraphs/tours/5}{It does have a Hamiltonian \emph{path}. Petersen
is the smallest \emph{hypohamiltonian} graph.}
\pair{graphs-hypergraphs/tours/6}{The \(3\times3\) grid has no Hamiltonian cycle
either, for a parity reason explained below.}
\pair{graphs-hypergraphs/tours/7}{\B{FindPath} with \mcode{All} lists every
simple path between two vertices: both ways round the pentagon.}
\end{codepairs}

The Gray code on the cube, highlighted:

\plotcell{graphs-hypergraphs/hamilton}


\calloutnote{Theory}{The grid is bipartite, and its two colour classes have
\(5\) and \(4\) vertices. A cycle in a bipartite graph alternates between the
classes, so a Hamiltonian cycle would need the two classes to be the same
size. No search is needed. The Petersen case has no such short proof. Its
non-Hamiltonicity is a classical result that is usually proved by examining
cases, which is what the search does.}

\calloutnote{Under the hood}{\B{FindHamiltonianCycle}
(\texttt{src/graph/galg\_hamilton.c}) is a depth-first search with
\emph{constraint propagation}. Each step decides whether one \emph{edge} is in
the cycle. A vertex with only two remaining edges forces both of them into the
cycle. A vertex that already has two chosen edges excludes all its others. An
edge that would close a cycle shorter than \(n\) is refused. At every node of
the search, the edges not yet excluded must form a biconnected graph, which
Tarjan's algorithm checks in linear time. The
search works under a node budget. If the budget runs out, the call returns
\emph{unevaluated}, so an empty list always means the search proved there is no
Hamiltonian cycle.}

\section{Centrality and clustering}
\label{sec:graphs-centrality}

A \emph{centrality}\index{centrality} assigns each vertex a number that measures
how important it is. There are many definitions, and they disagree in
informative ways. A small graph shows the differences: a triangle
\(\{1,2,3\}\) with a pendant vertex \(4\) attached to vertex \(3\).

\begin{codepairs}
\pair{graphs-hypergraphs/centrality/1}{A triangle with a tail.}
\pair{graphs-hypergraphs/centrality/2}{\B{DegreeCentrality} counts
neighbours.}
\pair{graphs-hypergraphs/centrality/3}{\B{ClosenessCentrality} is the reciprocal
of the mean distance to the other vertices. Vertex \(3\) is at distance \(1\)
from all of them.}
\pair{graphs-hypergraphs/centrality/4}{\B{BetweennessCentrality} counts shortest
paths through a vertex. Only vertex \(3\) lies between others: on the paths
\(1\)--\(4\) and \(2\)--\(4\).}
\pair{graphs-hypergraphs/centrality/5}{\B{PageRankCentrality} is the stationary
distribution of a random surfer. It sums to \(1\).}
\pair{graphs-hypergraphs/centrality/6}{\B{EigenvectorCentrality} is the Perron
vector of \(A\): a vertex is important if its neighbours are.}
\end{codepairs}

Every measure ranks vertex \(3\) first and vertex \(4\) last. They differ in
how large the gaps are. Betweenness gives every vertex except \(3\) a score of
zero, while PageRank still gives the leaf a noticeable share. \B{KatzCentrality}
and \B{HITSCentrality} (hubs and authorities, for directed graphs) complete the
set.

\emph{Clustering} measures how close each neighbourhood is to a clique. The
local coefficient of a vertex is the fraction of pairs of its neighbours that
are themselves adjacent:\index{clustering coefficient}

\begin{codepairs}
\pair{graphs-hypergraphs/centrality/7}{Vertices \(1\) and \(2\) have a complete
neighbourhood. Of the three pairs among vertex 3's neighbours, only one is an
edge. A leaf has no pairs.}
\pair{graphs-hypergraphs/centrality/8}{The global coefficient (three times the
triangles, divided by the connected triples) differs from the mean of the local
coefficients. Both are exact.}
\end{codepairs}

\calloutnote{Pitfall}{The \emph{global} and \emph{mean} clustering coefficients
answer different questions and are easy to confuse. The global coefficient
weights each vertex by the number of neighbour pairs it has, so hubs dominate.
The mean weights every vertex equally, so the many leaves of a sparse network
pull it down. Here the two are \(3/5\) and \(7/12\). On a real social network
they can differ by a factor of two.}

\calloutnote{Under the hood}{\B{BetweennessCentrality} is Brandes's algorithm:
one breadth-first search per source, followed by a reverse accumulation of
dependencies. It costs \(O(nm)\) in total. The sources are independent, so
Mathilda divides them among threads and adds up the per-thread vectors.
\B{EigenvectorCentrality} uses Arnoldi iteration on the sparse adjacency
structure and never builds the dense matrix. On a \(10^4\)-vertex graph,
betweenness takes \(0.23\) s against Mathematica's \(4.5\) s, and Katz
centrality on \(10^5\) vertices takes \(31\) ms against \(38\) s.}

\section{Random graphs and the giant component}
\label{sec:graphs-random}

The Erd\H{o}s--R\'enyi random graph \(G(n, m)\) chooses \(m\) of the
\(\binom{n}{2}\) possible edges uniformly at random. Its most famous property is
a \emph{phase transition}.\index{random graph!giant component} When the average
degree \(c = 2m/n\) is below \(1\), every connected component is small, of size
\(O(\log n)\). When \(c\) passes \(1\), a single \emph{giant component}
containing a positive fraction of all the vertices appears suddenly. A table of
the largest component in \(G(1000, m)\) shows the transition:

\begin{codepairs}
\pair{graphs-hypergraphs/random/1}{Largest component size as \(m\) grows. The
average degree reaches \(1\) at \(m = 500\).}
\end{codepairs}

A picture shows the same transition at a smaller size. With \(120\) vertices and
\(80\) edges the average degree is \(4/3\), just above the threshold. One
component of \(47\) vertices, drawn in red, has already absorbed a large share
of the graph, while the rest are small trees and isolated vertices:

\plotcell{graphs-hypergraphs/giant}


Below \(m = 500\) the largest component has about a dozen vertices. By
\(m = 600\) (average degree \(1.2\)) it holds a third of the graph, and by
\(m = 1000\) (average degree \(2\)) more than four-fifths. The theory predicts
the size. The fraction \(\beta\) of vertices in the giant component solves
\(\beta = 1 - e^{-c\beta}\), which gives \(\beta \approx 0.797\) at \(c = 2\).
That is close to the \(815\) that the session found.

The degree of a single vertex in \(G(n,m)\) is approximately Poisson with mean
\(c\). A sample with \(n = 200\) and \(c = 4\) is already close:

\begin{codepairs}
\pair{graphs-hypergraphs/random/2}{A seeded random graph with \(200\) vertices
and \(400\) edges.}
\pair{graphs-hypergraphs/random/3}{Its mean degree is exactly \(2m/n = 4\).}
\pair{graphs-hypergraphs/random/4}{The degree distribution: how many vertices
have each degree.}
\pair{graphs-hypergraphs/random/5}{The counts a Poisson distribution with mean
\(4\) predicts for \(200\) vertices.}
\end{codepairs}

\calloutnote{Performance}{\B{RandomGraph} samples \(m\) indices out of the
\(\binom{n}{2}\) candidate edges without replacement, making exactly the draws
\B{RandomSample} would make, and decodes each index to its vertex pair with a
binary search over the rows of the upper triangle. The candidate list is never
built, so memory is \(O(m)\), not \(O(n^2)\). A seeded \B{RandomGraph} gives the
same graph as sampling the explicit candidate list would.}

\section{Flows, cuts, and connectivity}
\label{sec:graphs-flows}

Think of each directed edge as a pipe with a \emph{capacity}. How much can flow
from a source \(s\) to a sink \(t\), if every intermediate vertex passes on
exactly what it receives? \B{FindMaximumFlow} takes the capacities as an
option, listed in \B{EdgeList} order:\index{maximum flow}

\begin{codepairs}
\pair{graphs-hypergraphs/flows/1}{A four-vertex network. The capacities are
supplied to each call below.}
\pair{graphs-hypergraphs/flows/2}{The maximum flow from \(s\) to \(t\) is
\(5\).}
\pair{graphs-hypergraphs/flows/3}{\mcode{"FlowMatrix"} shows how: both edges out
of \(s\) are full (\(3\) and \(2\)), and vertex \(a\) splits its \(3\) as \(1\)
to \(b\) and \(2\) to \(t\).}
\pair{graphs-hypergraphs/flows/4}{\B{FindEdgeCut} gives the fewest edges whose
removal separates \(s\) from \(t\). Here those are the two edges out of
\(s\).}
\end{codepairs}

Drawn in layers from source to sink, with each edge labelled by its flow and
its capacity, the saturated cut is easy to see. Both edges out of \(s\) carry
their full capacity:

\plotcell{graphs-hypergraphs/flow}


The flow value is \(5\), and the two edges out of \(s\) have total capacity
\(3 + 2 = 5\). This is not a coincidence.

\begin{theorem}[Max-flow min-cut]
\index{max-flow min-cut theorem}
In any network, the value of a maximum \(s\)--\(t\) flow equals the minimum total
capacity of a set of edges whose removal separates \(s\) from \(t\).
\end{theorem}

Applied to unit capacities, the theorem becomes \emph{Menger's theorem}. The
number of edge-disjoint paths between two vertices equals the minimum number of
edges whose removal separates them. This turns connectivity into a flow
computation:

\begin{codepairs}
\pair{graphs-hypergraphs/flows/5}{The Petersen graph again.}
\pair{graphs-hypergraphs/flows/6}{\B{FindMinimumCut} is a global minimum cut: the
cheapest way to split the graph in two. Isolating one vertex costs its three
edges.}
\pair{graphs-hypergraphs/flows/7}{\B{FindVertexCut} gives the three neighbours
of a vertex. Removing them disconnects that vertex from the rest.}
\pair{graphs-hypergraphs/flows/8}{\B{VertexConnectivity}, \B{EdgeConnectivity},
and the minimum degree are all \(3\).}
\end{codepairs}

Whitney's inequality \(\kappa(G) \le \lambda(G) \le \delta(G)\) holds for every
graph: vertex connectivity is at most edge connectivity, which is at most the
minimum degree. The Petersen graph attains equality throughout. It is as well
connected as a 3-regular graph can be, which is part of why it appears so often
in network design.

\calloutnote{Under the hood}{\B{FindMaximumFlow} is \emph{Dinic's algorithm}. It
builds breadth-first level graphs, truncated at the sink, and pushes blocking
flows along them using current-arc pointers on a compressed residual network.
Every capacity is converted to a 64-bit integer first. Rational and real
capacities are scaled by a common factor, so even the maximum flow of
floating-point capacities is computed without rounding error.
\B{FindMinimumCut} uses the Nagamochi--Ibaraki algorithm, which needs no source
or sink. \B{VertexConnectivity} uses Even's reduction to \(O(n)\) flow problems
and replaces an earlier search over vertex subsets whose cost grew
exponentially.}

\section{Matchings, covers, and cliques}
\label{sec:graphs-covers}

A \emph{matching}\index{matching} is a set of edges that share no vertex. A
\emph{vertex cover}\index{vertex cover} is a set of vertices that touches every
edge. An \emph{independent set}\index{independent set} is a set of vertices no
two of which are adjacent. A \emph{clique} is a set of vertices every two of
which are adjacent. Maximum matchings can be found in polynomial time. The other
three problems are NP-hard. Mathilda solves all four exactly:

\begin{codepairs}
\pair{graphs-hypergraphs/covers/1}{Petersen once more.}
\pair{graphs-hypergraphs/covers/2}{\B{FindIndependentEdgeSet} gives a maximum
matching. Five edges cover all ten vertices, so the matching is
\emph{perfect}.}
\pair{graphs-hypergraphs/covers/3}{\B{FindVertexCover} gives a minimum vertex
cover of six vertices\ldots}
\pair{graphs-hypergraphs/covers/4}{\ldots and \B{FindIndependentVertexSet} gives
a maximum independent set of four.}
\pair{graphs-hypergraphs/covers/5}{Together they make up the whole vertex
set.}
\pair{graphs-hypergraphs/covers/6}{\B{FindClique} finds the largest clique. In
the Petersen graph that is a single edge: the graph has no triangles.}
\pair{graphs-hypergraphs/covers/7}{In \(K_5\) the whole graph is a clique.}
\pair{graphs-hypergraphs/covers/8}{The chromatic number is \(3\), because the
graph contains 5-cycles, which need three colours.}
\end{codepairs}

On the classic drawing the three structures are easy to compare. The minimum
vertex cover is in red, and the four blue vertices left over are the maximum
independent set:

\plotcell{graphs-hypergraphs/cover}

The perfect matching pairs off all ten vertices with five disjoint edges:

\plotcell{graphs-hypergraphs/matching}

And a proper colouring with three colours, one per class of the partition that
\B{FindVertexColoring} returns:

\plotcell{graphs-hypergraphs/coloring}


Line 5 is Gallai's identity. A set is independent exactly when its complement is
a vertex cover. So a maximum independent set and a minimum cover are
complements of each other, and \(\alpha(G) + \tau(G) = n\), here \(4 + 6 = 10\).
A clique in \(G\) is an independent set in the complement of \(G\), so the four
problems reduce to two.

\calloutnote{Theory}{For \emph{bipartite} graphs the NP-hard problems become
easy. \emph{K\H{o}nig's theorem} says that the minimum vertex cover of a
bipartite graph has the same size as the maximum matching, and a maximum
matching can be found in polynomial time. Mathilda uses this. When the
independent-set solver finds a bipartite connected component of \(256\) or more
vertices, it solves that component with Hopcroft--Karp and K\H{o}nig's
construction instead of branching on it.}

\calloutnote{Under the hood}{Matching uses Hopcroft--Karp for bipartite graphs
and Edmonds's blossom algorithm otherwise. Independent set and vertex cover use
\emph{branch and reduce}: degree-0, degree-1, and degree-2 folding, domination,
and connected-component splitting, with the search itself kept for what remains.
\B{FindClique} uses BBMC, a bitset branch and bound that colours the candidate
set to prune. Every solver has a budget and returns unevaluated when it runs
out, never a smaller answer passed off as optimal. A randomised differential
test of about 12\,000 cases against Mathematica~15 found no disagreements in
optimum sizes. It did find that Mathematica sometimes returns non-maximum
matchings on weighted graphs.}

\section{Isomorphism, symmetry, and planarity}
\label{sec:graphs-iso}

Two graphs are \emph{isomorphic}\index{graph isomorphism} if some relabelling of
the vertices turns one into the other. Vertex names and the order in which
edges are listed do not affect the answer. \B{IsomorphicGraphQ} decides the
question, \B{FindGraphIsomorphism} returns the relabelling, and
\B{CanonicalGraph} gives a canonical representative. Two graphs are isomorphic
exactly when their canonical forms are identical:

\begin{codepairs}
\pair{graphs-hypergraphs/isomorphism/1}{A square on symbolic vertices, with its
edges listed out of order.}
\pair{graphs-hypergraphs/isomorphism/2}{It is the 4-cycle.}
\pair{graphs-hypergraphs/isomorphism/3}{The isomorphism is returned as an
\B{Association} from old vertex names to new ones.}
\pair{graphs-hypergraphs/isomorphism/4}{The two canonical forms are identical.
\B{SameQ} is enough to compare them.}
\pair{graphs-hypergraphs/isomorphism/5}{\B{GraphAutomorphismGroup} gives the
symmetries of the Petersen graph as a permutation group with four
generators.}
\end{codepairs}

That group has order \(120\) and is isomorphic to the symmetric group \(S_5\).
The Petersen graph is the Kneser graph \(K(5,2)\): its vertices are the
2-subsets of a 5-set, and two vertices are adjacent when the subsets are
disjoint. Any permutation of the five points gives a symmetry.

\calloutnote{Under the hood}{Isomorphism, canonical form, and automorphisms all
use one \emph{individualization--refinement} engine, the method used by nauty
and bliss. Colour refinement (1-dimensional Weisfeiler--Leman) splits the
vertices by degree profile until the partition is stable. If some colour class
still has more than one vertex, the engine singles out one vertex, refines
again, and searches the resulting tree, pruning with the automorphisms found so
far. On random 3-regular graphs with \(10^4\) vertices, which refinement alone
cannot fully separate, Mathilda takes \(4\) ms and Mathematica \(818\) ms.}

A graph is \emph{planar}\index{planar graph} if it can be drawn in the plane
without edge crossings. By Kuratowski's theorem, a graph is planar exactly when
it contains no subdivision of \(K_5\) or of \(K_{3,3}\):

\begin{codepairs}
\pair{graphs-hypergraphs/isomorphism/6}{\(K_4\) and the grid are planar. The two
Kuratowski graphs \(K_5\) and \(K_{3,3}\) are not.}
\pair{graphs-hypergraphs/isomorphism/7}{The Petersen graph is not planar. It
contains a subdivision of \(K_{3,3}\).}
\end{codepairs}

The two Kuratowski graphs, drawn by \B{GraphPlot}. However the vertices of
\(K_5\) are placed, some pair of edges must cross:

\plotcell{graphs-hypergraphs/k5}

\(K_{3,3}\) in the two-column bipartite layout. It is the ``three utilities''
puzzle: three houses, three utilities, and no way to connect each house to each
utility without a crossing:

\plotcell{graphs-hypergraphs/k33}


\calloutnote{Under the hood}{\B{PlanarGraphQ} is the \emph{left--right planarity
test} of de~Fraysseix and Rosenstiehl, as refined by Brandes. It runs in linear
time, in two depth-first passes: one orients the edges and computes low-points,
and the other checks that the back edges can be split consistently between the
left and right sides of each tree edge. It never builds an embedding, so it
never needs to handle one.}

\section{Graphs are expressions}
\label{sec:graphs-programming}

Everything in this chapter so far used graph functions. But a graph is an
ordinary expression, so the rest of the language works on it too. Graphs can be
built from data with \B{KeyValueMap} and \B{Thread}, filtered with \B{Select},
and rewritten with rules:

\begin{codepairs}
\pair{graphs-hypergraphs/programming/1}{A small ``who follows whom'' table, held
as an association from each person to a list.}
\pair{graphs-hypergraphs/programming/2}{Thread each key over its list of values
to get the directed edges.}
\pair{graphs-hypergraphs/programming/3}{The edges, in the order the data gave
them.}
\pair{graphs-hypergraphs/programming/4}{People nobody follows: vertices of
in-degree zero.}
\pair{graphs-hypergraphs/programming/5}{\B{EdgeRules} goes the other way and
turns edges back into rules.}
\pair{graphs-hypergraphs/programming/6}{A rule rewrites edges like any other
expression.}
\pair{graphs-hypergraphs/programming/7}{\B{NeighborhoodGraph} is the ball of
radius \(1\) around vertex \(1\): the vertex and its three neighbours.}
\pair{graphs-hypergraphs/programming/8}{\B{Subgraph} keeps the chosen vertices
and the edges among them. In this labelling, vertices \(1\) to \(5\) form a
5-cycle.}
\end{codepairs}

\calloutnote{Under the hood}{A rule applied to an \B{EdgeList} produces a list of
edges, not a graph. Wrap the result in \B{Graph} to validate it again. This is
deliberate. \B{Graph} is the single point where simplicity and endpoint
membership are checked, and where the validated-graph memo is seeded. Every
function that accepts a graph can therefore assume the invariants hold, and does
not need to check them again.}

\section{Hypergraphs}
\label{sec:hypergraphs}

In a graph every edge has two ends. Many relations have more participants. A
paper has several authors, a reaction has several reagents, and a row of a
database relates several fields. Replacing such a relation with all the pairwise
edges it implies loses information: three authors on one paper and three
separate two-author papers produce the same triangle. A
\emph{hypergraph}\index{hypergraph} keeps the relation intact. It consists of a
vertex set and a list of \emph{hyperedges}, and each hyperedge is an arbitrary
subset of the vertices.

\subsection{Building and reading a hypergraph}
\label{ssec:hyp-construct}

\B{Hypergraph} takes a list of hyperedges, and optionally a vertex list first.
Like \B{Graph}, it is an ordinary expression, validated when it is built:

\begin{codepairs}
\pair{graphs-hypergraphs/hyper-basics/1}{Four hyperedges of sizes \(3, 2, 3, 1\)
on seven vertices.}
\pair{graphs-hypergraphs/hyper-basics/2}{The canonical form lists the vertices
first, as \B{Graph} does.}
\pair{graphs-hypergraphs/hyper-basics/3}{The graph accessors also work on
hypergraphs. \B{VertexDegree} counts the hyperedges containing each
vertex.}
\pair{graphs-hypergraphs/hyper-basics/4}{\B{HyperedgeSizes} gives the size of
each hyperedge.}
\pair{graphs-hypergraphs/hyper-basics/5}{\B{HypergraphRank} and
\B{HypergraphCorank} are the largest and smallest hyperedge sizes. They differ,
so the hypergraph is not \emph{uniform}.}
\pair{graphs-hypergraphs/hyper-basics/6}{\B{IncidenceMatrix} is \(7\times 4\):
entry \((i, j)\) is \(1\) when vertex \(i\) is in hyperedge \(j\).}
\pair{graphs-hypergraphs/hyper-basics/7}{\B{HypergraphDual} exchanges the roles
of vertices and hyperedges. Vertex 3 lies in hyperedges 1 and 2, so it becomes
the dual hyperedge \(\{1, 2\}\).}
\end{codepairs}

\mcode{HypergraphPlot} draws each hyperedge as a translucent region around its
members: a rounded hull for three or more vertices, a stadium for two, and a
ring for one. Vertices \(3\) and \(4\) each lie in two hyperedges, and vertex
\(7\) sits alone in a singleton:

\plotcell{graphs-hypergraphs/hyper}


The incidence matrix fully describes the hypergraph, and the dual has the
transposed incidence matrix. Applying the dual twice returns the original
hypergraph, with the vertices renamed by their positions. A hypergraph can equally be read as a family of sets or as a
bipartite graph between vertices and hyperedges, and each hypergraph operation
is easiest to understand in one of these two readings.

\calloutnote{Pitfall}{A hyperedge list is kept exactly as written, including its
order and any repeated vertex. This lets Wolfram-model states such as
\mcode{\{\{1,2,3\},\{3,4\}\}} be used directly. Most hypergraph operations treat
a hyperedge as the \emph{set} of its distinct vertices. The order-sensitive ones
are \B{HyperedgeSizes}, which counts repeats, and \B{HypergraphToGraph}, which
reads a hyperedge as an ordered chain.}

\subsection{From hypergraphs to graphs}
\label{ssec:hyp-expansions}

There are three standard ways to turn a hypergraph into a graph, and each keeps
different information. The \emph{clique expansion} makes every hyperedge a
clique. The \emph{star expansion} adds one vertex for each hyperedge and joins
it to that hyperedge's members, which is the bipartite incidence graph. The
Wolfram Function Repository convention in \B{HypergraphToGraph} reads each
hyperedge as an ordered chain:

\begin{codepairs}
\pair{graphs-hypergraphs/hyper-structure/1}{Name the hypergraph from the previous
section.}
\pair{graphs-hypergraphs/hyper-structure/2}{\B{HypergraphCliqueExpansion}: the
two triples become triangles and the pair becomes an edge. The singleton
\(\{7\}\) adds no edges.}
\pair{graphs-hypergraphs/hyper-structure/3}{\B{HypergraphStarExpansion} has
\(7 + 4\) vertices and one edge for each incidence, \(3+2+3+1 = 9\).}
\pair{graphs-hypergraphs/hyper-structure/4}{\B{HypergraphToGraph} turns an
ordered hyperedge \(\{a, b, c\}\) into the directed edges \(a\to b\),
\(a \to c\), \(b \to c\).}
\end{codepairs}

Connectivity and distance make sense without expanding. Two vertices are
adjacent when some hyperedge contains both, and two hyperedges are adjacent
when they intersect:

\begin{codepairs}
\pair{graphs-hypergraphs/hyper-structure/5}{\B{HypergraphConnectedComponents}:
the chain \(\{1,2,3\}\), \(\{3,4\}\), \(\{4,5,6\}\) connects six vertices.
Vertex \(7\) is isolated.}
\pair{graphs-hypergraphs/hyper-structure/6}{Vertex \(1\) to vertex \(6\) takes
three hyperedges. Hyperedge \(1\) to hyperedge \(3\) takes two steps, through
hyperedge \(2\).}
\end{codepairs}

\subsection{\texorpdfstring{\(s\)}{s}-line graphs}
\label{ssec:hyp-sline}

The \emph{line graph} of a hypergraph has one vertex for each hyperedge, and
two of them are joined when the hyperedges intersect. Asking for an overlap of
at least \(s\) vertices gives the \emph{\(s\)-line graph}\index{hypergraph!s-line
graph}, a family of increasingly strict views of the same data. For
co-authorship data, \(s = 1\) links papers that share any author and \(s = 2\)
links only papers that share a pair of authors. Consider four hyperedges
arranged around a ring:

\begin{codepairs}
\pair{graphs-hypergraphs/hyper-structure/7}{A ring of four hyperedges:
\(\{1,2,3\}\), \(\{2,3,4\}\), \(\{3,4,5\}\), \(\{1,5\}\).}
\pair{graphs-hypergraphs/hyper-structure/8}{\B{HypergraphLineGraph} with
\(s = 1\) has five of the six possible edges. Only hyperedges \(2\) and \(4\) are
disjoint.}
\pair{graphs-hypergraphs/hyper-structure/9}{With \(s = 2\) only overlaps of two
or more vertices count: \(\{2,3\}\) and \(\{3,4\}\).}
\pair{graphs-hypergraphs/hyper-structure/10}{\B{HyperedgeConnectedComponents} at
\(s = 2\): the hyperedge \(\{1,5\}\) is no longer connected to the others.}
\end{codepairs}

The ring hypergraph drawn. Hyperedges \(1\), \(2\), and \(3\) overlap in two
vertices each, which is exactly what the \(2\)-line graph detects:

\plotcell{graphs-hypergraphs/ring}


\calloutnote{Performance}{Most hypergraph operations run in time linear in the
total incidence \(\sum_e |e|\). The \(s\)-overlap family is the exception: it
costs \(O(\sum_v \deg(v)^2)\), the number of hyperedge pairs that meet at some
vertex. The \(1\)-line graph has that many edges anyway, so no algorithm could
do better. The algorithm counts overlaps for each vertex's incidence list in a
scratch array and never forms the intersection of two hyperedges directly.}

\subsection{Transversals}
\label{ssec:hyp-transversals}

A \emph{transversal}\index{hypergraph!transversal} (or \emph{hitting set}) is a
set of vertices that meets every hyperedge. The minimal transversals form a
hypergraph of their own, the \emph{transversal hypergraph}
\(\operatorname{Tr}(H)\). This single construction covers many problems. The
minimal vertex covers of a graph are the transversals of its edge set. The
minimal diagnoses of a faulty circuit are the transversals of its conflict sets.
In a database, the candidate keys and the minimal non-keys are transversal
hypergraphs of each other.

\begin{codepairs}
\pair{graphs-hypergraphs/hyper-transversals/1}{The ring hypergraph from above.}
\pair{graphs-hypergraphs/hyper-transversals/2}{\B{TransversalHypergraph}: the
four minimal hitting sets. Each of them meets all four hyperedges.}
\pair{graphs-hypergraphs/hyper-transversals/3}{\B{FindMinimumTransversal} gives
a smallest hitting set. Its size, \(2\), is the transversal number
\(\tau(H)\).}
\pair{graphs-hypergraphs/hyper-transversals/4}{Taking transversals twice gives
back the original hyperedges, as Berge's duality theorem says it must.}
\end{codepairs}

\begin{theorem}[Berge]
If \(H\) is a \emph{simple} hypergraph (no hyperedge contains another), then
\(\operatorname{Tr}(\operatorname{Tr}(H)) = H\).
\end{theorem}

The ring is simple, and the last line of the session confirms the theorem.
Computing \(\operatorname{Tr}(H)\) is the \emph{hypergraph dualization} problem.
Its exact complexity is a well-known open question. Fredman and Khachiyan showed
that it can be solved in quasi-polynomial time in the combined size of the input
and output, and no polynomial algorithm is known.

\calloutnote{Under the hood}{\B{TransversalHypergraph} is \emph{MMCS}
(minimal-to-minimal via crit and uncov, Murakami and Uno). It grows a candidate
set one vertex at a time, and for each chosen vertex it keeps track of the
hyperedges only that vertex hits. A vertex that stops having any such hyperedge
cannot belong to a minimal transversal, so that branch is cut immediately. On a
16-vertex instance it runs in \(4.2\) ms. The Function Repository's
\mcode{TransversalHypergraph} takes \(5.2\) s. \B{FindMinimumTransversal} is a
combinatorial branch and bound. On an \(80\)-vertex instance it runs
\(3.2\)--\(3.6\times\) faster than the MILP solvers (Mathematica's
\mcode{LinearOptimization}, SciPy's HiGHS) that are the usual way to compute
it.}

Random hypergraphs and sub-hypergraphs complete the toolkit.
\B{RandomHypergraph}\mcode{[\{n, m\}, k]} samples a \(k\)-uniform hypergraph and
uses the stream that \B{SeedRandom} seeds. There are two ways to restrict a
hypergraph to a vertex subset, and they are easy to confuse:

\begin{codepairs}
\pair{graphs-hypergraphs/hyper-transversals/5}{A reproducible random
3-uniform hypergraph: eight vertices, five triples.}
\pair{graphs-hypergraphs/hyper-transversals/6}{Its hyperedges.}
\pair{graphs-hypergraphs/hyper-transversals/7}{\B{Subhypergraph} keeps only the
hyperedges that lie \emph{entirely} inside the subset\ldots}
\pair{graphs-hypergraphs/hyper-transversals/8}{\ldots while
\B{HypergraphRestriction} intersects \emph{every} hyperedge with the subset and
drops only the ones that become empty.}
\end{codepairs}

\calloutnote{Theory}{Mathematica~15 has no built-in hypergraph type:
\mcode{Names["*Hypergraph*"]} returns an empty list. Hypergraphs are available
in the Wolfram ecosystem only as top-level add-ons: the
WolframInstitute/Hypergraph paclet, Function Repository entries such as
\mcode{ResourceFunction["TransversalHypergraph"]}, and the Wolfram Physics
Project tools. Mathilda uses their names where one exists, accepts their
argument forms, and implements everything in C.}

\section{Featured examples}
\label{sec:graphs-featured}

The examples in this section are short. Each is a small, self-contained problem
from outside mathematics, solved in a few lines with the functions of this
chapter.

\subsection*{The knight's tour}
\index{knight's tour}

Can a knight visit every square of a chessboard exactly once and return to its
starting square? Build the graph whose vertices are the \(64\) squares, with an
edge for every knight move, and the question becomes: is this graph Hamiltonian?

\begin{codepairs}
\pair{graphs-hypergraphs/featured/1}{The four ``forward'' knight moves. Each edge
is generated once.}
\pair{graphs-hypergraphs/featured/2}{The knight's graph of an \(n\times n\)
board, with squares named \(\{i, j\}\).}
\pair{graphs-hypergraphs/featured/3}{The standard board has \(168\) knight
moves.}
\pair{graphs-hypergraphs/featured/4}{Corner squares have only two moves, and
central squares have eight.}
\pair{graphs-hypergraphs/featured/5}{A closed tour exists: a Hamiltonian cycle
through all \(64\) squares.}
\pair{graphs-hypergraphs/featured/6}{On the \(5\times5\) board there is no closed
tour, but an open one visits all \(25\) squares.}
\end{codepairs}

Placing each square at its own board coordinates turns the graph drawing into a
chessboard. The grey lines are all \(168\) knight moves, and the red ones are
the closed tour that \B{FindHamiltonianCycle} found:

\plotcell{graphs-hypergraphs/knight}


The \(5 \times 5\) result has a one-line explanation. A knight always moves to a
square of the other colour, so the knight's graph is bipartite. A closed tour
alternates colours, so it needs as many light squares as dark ones. A board with
\(25\) squares cannot have that.

\subsection*{Find the powerbrokers of Renaissance Florence}
\index{Florentine families}

In a well-known study of fifteenth-century Florence, Padgett and Ansell recorded
the marriages between the city's leading families. Betweenness centrality
explains how the Medici came to dominate:

\begin{codepairs}
\pair{graphs-hypergraphs/featured/7}{Twenty marriage ties among fifteen
families.}
\pair{graphs-hypergraphs/featured/8}{The Medici lie on far more shortest paths
than any other family: \(47.5\), twice the next highest.}
\pair{graphs-hypergraphs/featured/9}{By degree alone the Medici lead less
dramatically.}
\pair{graphs-hypergraphs/featured/10}{Removing the Medici alone disconnects the
network.}
\pair{graphs-hypergraphs/featured/11}{A marriage alliance between the Pazzi and
the Lamberteschi has to pass through the Medici.}
\end{codepairs}

The marriage network, with the Medici in red. Every path from the Pazzi or the
Salviati to the rest of Florence passes through them:

\plotcell{graphs-hypergraphs/florence}


The Medici had only a few more marriages than the Strozzi or the Guadagni.
What set them apart was their position: they were the family through which
everyone else's connections ran, and that is what betweenness measures.

\subsection*{Colour the map of Australia}

Colour Australia's states and territories so that neighbours differ. The
constraints form a graph with one vertex per region. Tasmania is an island, so
it has no neighbours.

\begin{codepairs}
\pair{graphs-hypergraphs/featured/12}{Seven regions and their land borders.}
\pair{graphs-hypergraphs/featured/13}{Three colours suffice, and three are
needed: WA, NT, and SA form a triangle.}
\end{codepairs}

With each region placed roughly where it is on the map and coloured by
\B{FindVertexColoring}:

\plotcell{graphs-hypergraphs/australia}


\subsection*{Compute collaboration distance}

A mathematician's \emph{Erd\H{o}s number} is their collaboration distance from
Paul Erd\H{o}s. Papers are hyperedges over their authors, and collaboration
distance is hypergraph distance:

\begin{codepairs}
\pair{graphs-hypergraphs/featured/14}{Five papers as hyperedges over seven
authors.}
\pair{graphs-hypergraphs/featured/15}{From R\'enyi to Patashnik takes three
papers.}
\pair{graphs-hypergraphs/featured/16}{A smallest set of authors that covers
every paper. No two authors cover all five.}
\end{codepairs}

The five papers drawn as regions. Graham sits in three of them, which is why
every short collaboration path passes near him:

\plotcell{graphs-hypergraphs/papers}


\subsection*{What a graph forgets}

Why use hypergraphs at all? Because the pairwise view loses information. One
paper with three authors and three separate two-author papers produce the same
collaboration graph:

\begin{codepairs}
\pair{graphs-hypergraphs/featured/17}{The two clique expansions are isomorphic,
although the hypergraphs are different.}
\end{codepairs}

\subsection*{Get dressed in the right order}

The dependency graph of \S\ref{sec:graphs-directed} is itself a featured
example. \B{TopologicalSort} returns an order that respects every constraint,
and adding one contradictory constraint makes \B{AcyclicGraphQ} report that no
order exists.

\section{How fast is it?}
\label{sec:graphs-performance}

Every head in this chapter was benchmarked against Mathematica~15.0 and against
the standard Python libraries (networkx for graphs, xgi for hypergraphs). The
benchmark sources are in \texttt{benchmarks/92-graph-predicates} through
\texttt{benchmarks/96-hypergraphs}, and each row was checked for agreement of
values across the systems. Table~\ref{tab:graph-bench} gives a representative
selection. Times are milliseconds on Apple silicon, best of three.

\begin{table}[ht]
\centering
\small
\begin{tabular}{@{}lrrr@{}}
\toprule
Case & Mathilda & Mathematica & ratio \\
\midrule
\B{VertexDelete}, \(1000\) of \(10^5\) vertices & 4 & 52 & \(13\times\) \\
\B{GraphUnion}, two \(10^5\)-edge graphs & 13 & 354 & \(27\times\) \\
\B{BetweennessCentrality}, \(10^4\) vertices & 230 & 4\,500 & \(20\times\) \\
\B{KatzCentrality}, \(10^5\) vertices & 31 & 38\,000 & \(1200\times\) \\
\B{IsomorphicGraphQ}, 3-regular, \(10^4\) vertices & 4 & 818 & \(200\times\) \\
Maximum matching, triangulated grid & 2.8 & 190 & \(68\times\) \\
\midrule
\B{Hypergraph} construction, \(10^5\) hyperedges & 19.9 & 2\,912 & \(146\times\) \\
\B{HypergraphDual}, \(10^5\) hyperedges & 8.1 & 236 & \(29\times\) \\
\B{HyperedgeDistance}, \(10^5\) hyperedges & 2.8 & 7\,255 & \(2600\times\) \\
\B{TransversalHypergraph}, 16 vertices & 4.2 & 5\,219 & \(1200\times\) \\
\B{RandomHypergraph} (FR form), \(10^5\) triples & 10.6 & 1.6 & \(0.15\times\) \\
\bottomrule
\end{tabular}
\caption{Representative graph and hypergraph timings, in milliseconds. The
hypergraph baselines are the Wolfram paclet and Function Repository code, which
is top-level Wolfram Language, not kernel C.}
\label{tab:graph-bench}
\end{table}

Two caveats keep the table fair. First, the very large hypergraph ratios compare
compiled C with interpreted Wolfram Language. That is the real choice facing
Mathematica users, but it is not a comparison of equally engineered kernels. The
graph rows, by contrast, compare with Mathematica's own C++ kernel, and those are
the more meaningful figures. Second, Mathilda loses one row. The Function
Repository's \mcode{RandomHypergraph} returns a single packed integer array.
Mathilda returns a validated \B{Hypergraph} whose \(10^5\) hyperedges are
separate expression nodes, and allocating those nodes alone takes longer than
\(1.6\) ms. \B{GraphComplement} is at parity. Every other benchmarked case is
at or ahead of both Mathematica and Python, both warm and cold (with a freshly
built graph before each timed call).

\calloutnote{Performance}{Most of the speed comes from representation, not from
clever algorithms. A graph stays an ordinary expression, so the
uniform-expression model of Chapter~\ref{ch:datastructures} is kept. But the
first operation on a graph resolves it into integer arrays in the
validated-graph memo, and every algorithm afterwards runs on those arrays with
no hashing and no pattern matching. The hypergraph memo in
\texttt{src/graph/hyp\_util.c} works the same way, and when the vertices are a
compact range of machine integers it looks them up by direct addressing.}

\subsection*{Where this connects}

This chapter builds on several earlier ones. The Laplacian is an ordinary
integer matrix, so the exact \B{Det} and \B{Eigenvalues} of \S\ref{sec:linalg}
gave the matrix-tree theorem directly. Isomorphisms come back as associations
(Chapter~\ref{ch:datastructures}), and random graphs and hypergraphs use the
same seeded stream as the rest of the system. The chapter also followed the
rule the rest of the book follows: the exact solvers for Hamiltonian cycles,
cliques, independent sets, covers, and transversals return an exact answer or
return unevaluated, and never return a result that only looks plausible. To draw
any of these graphs, see Chapter~\ref{ch:graphics}. For the options of any
function mentioned here, type \texttt{?Name} at the prompt to read its
reference page.
